%=====================================================================
% Paper Preamble: Optimized for Cryptography Research
%=====================================================================

%----- 0. 错误屏蔽与性能优化 -----%
\usepackage{silence}
\WarningFilter{latex}{Invalid UTF-8 byte}
\WarningFilter{latex}{`h' float specifier changed to `ht'}

%----- 1. 语言控制逻辑 -----%
\usepackage{etoolbox}
\providecommand{\ifChinese}{\iffalse} % 兜底定义，防止主文件没写 \newif

\ifChinese
    \usepackage[UTF8, punct]{ctex}
\fi

%----- 2. 数学包 (Types: PascalCase) -----%
\usepackage{amsmath, amsthm, amsfonts, amssymb}
\usepackage{mathrsfs, amsbsy, mathtools}
\usepackage{autobreak}
\usepackage{siunitx}

%----- 3. 图片与颜色 -----%
\usepackage{graphicx}
\usepackage{subcaption} 
\usepackage{xcolor}
\usepackage{float}

%----- 4. 表格增强 -----%
\usepackage{caption}
\usepackage{colortbl, makecell, tabularx, multirow, threeparttable, booktabs}

%----- 5. 标注、列表与排版 -----%
\usepackage{soul}
\usepackage[normalem]{ulem}
\usepackage{pifont} % 必须加载，用于对勾和叉
\usepackage{fullpage, setspace, framed, mdframed}
\usepackage{enumitem}
\setlist{nosep}

%----- 6. 算法环境 -----%
\usepackage[ruled, vlined, linesnumbered, algo2e]{algorithm2e}
\ifChinese
    \SetAlgorithmName{算法}{算法}{算法列表}
\else
    \SetAlgorithmName{Algorithm}{Algorithm}{List of Algorithms}
\fi

%----- 7. TikZ 绘图 -----%
\usepackage{tikz, pgfplots}
\pgfplotsset{compat=1.18}
\usetikzlibrary{
    matrix, fit, calc, patterns, backgrounds, 
    arrows.meta, shapes, chains, decorations.pathreplacing, 
    intersections, positioning, external
}

\tikzset{
    >=Stealth,
    shapenode/.style = {draw, rectangle, fill=none, minimum size=0.5cm, minimum height=0.6cm, minimum width=0.6cm, 
    auto, node distance=0em, font=\normalsize}, 
    roundnode/.style = {draw, rectangle, rounded corners=0.5em, inner sep = 0em, 
      node distance=0em, minimum height=1.5em},
    rectanglenode/.style = {draw, rectangle, inner sep = 0em, node distance=0em},
    ellipsenode/.style = {draw, ellipse, inner sep = 0em, node distance=0em},
    textnode/.style  = {draw=none, fill=none, rectangle, node distance=0em, font=\normalsize},
    smalltextnode/.style  = {draw=none, fill=none, rectangle, minimum size=0.5cm, auto, node distance=0em, font=\small},
    circlenode/.style = {draw, circle, auto, inner sep = 0em, node distance=0em, minimum size=1.5em}, 
    dotnode/.style = {draw, fill=black, circle, auto, minimum size=0.1em, inner sep = 0em, node distance=0em},  
    connect/.style = {black,->}, 
    every text node part/.style={align=center}
}

%----- 8. 定理环境 -----%
\ifChinese
    \newtheorem{theorem}{定理}[section]
    \newtheorem{lemma}[theorem]{引理}
    \newtheorem{claim}[theorem]{断言}
    \newtheorem{proposition}[theorem]{命题}
    \newtheorem{corollary}[theorem]{推论}
    \newtheorem{definition}{定义}[section]
    \newtheorem{example}{例}[section]
    \newtheorem{remark}{注}[section]
    \renewcommand{\proofname}{\textbf{证明}}
    \newenvironment{solution}{\begin{proof}[\textbf{解}]}{\end{proof}}
\else
    \theoremstyle{plain}
    \newtheorem{theorem}{Theorem}[section]
    \newtheorem{lemma}[theorem]{Lemma}
    \newtheorem{claim}[theorem]{Claim}
    \newtheorem{proposition}[theorem]{Proposition}
    \newtheorem{corollary}[theorem]{Corollary}
    \theoremstyle{definition}
    \newtheorem{definition}{Definition}[section]
    \newtheorem{example}{Example}[section]
    \theoremstyle{remark}
    \newtheorem{remark}{Remark}[section]
\fi

%----- 9. 符号补丁 (修复日志中的 Undefined Control Sequence) -----%
% 修复 \checkmark 和 \xmark 的定义，确保表格正常
\let\oldcheckmark\checkmark
\renewcommand{\checkmark}{\text{\ding{51}}}
\newcommand{\xmark}{\text{\ding{55}}}

%----- 10. 颜色与密码学宏 -----%
\definecolor{LightMidnightBlue}{rgb}{0.2, 0.2, 0.7}
\definecolor{DarkGreen}{rgb}{0.0,0.5,0.0}

\newcommand{\Oracle}{\mathcal{O}}
\newcommand{\blue}[1]{{\color{blue}#1}}
\newcommand{\red}[1]{{\color{red}#1}}
\newcommand{\kP}{\mathbf{P}}
\newcommand{\kNP}{\mathbf{NP}}
\newcommand{\define}{\stackrel{\textup{def}}{=}}
\newcommand{\sample}{\xleftarrow{\textup{\tiny R}}}
\newcommand{\poly}{\textup{poly}}
\newcommand{\AdvA}{\mathsf{Adv}_\mathcal{A}}
\newcommand{\TaihangAssert}[1]{\texttt{TAIHANG\_ASSERT}(#1)}

%----- 11. 引用与超链接 (加载顺序关键) -----%
\usepackage[numbers, sectionbib, sort&compress]{natbib}
\setlength{\bibsep}{0.5ex}

\makeatletter
    \renewcommand\@biblabel[1]{{[#1]\hfill}}
\makeatother

\usepackage[colorlinks, anchorcolor=blue, linkcolor=blue, citecolor=blue, urlcolor=blue]{hyperref}
\usepackage[title, titletoc, header]{appendix}

\usepackage[capitalize]{cleveref}
\ifChinese
    \crefname{theorem}{定理}{定理}
    \crefname{figure}{图}{图}
    \crefname{table}{表}{表}
    \crefname{algorithm}{算法}{算法}
\fi

%----- 12. 自定义环境 -----%
\newenvironment{MyItemize}{\begin{itemize}[leftmargin=*, itemsep=1pt, parsep=0pt]}{\end{itemize}}
\newenvironment{MyEnumerate}{\begin{enumerate}[leftmargin=*, itemsep=1pt, parsep=0pt]}{\end{enumerate}}

\makeatletter
\@ifclassloaded{beamer}{}{
    \newenvironment{noteblock}{
        \begin{mdframed}[backgroundcolor=gray!10, linecolor=gray!50, roundcorner=2pt, innerleftmargin=5pt, innerrightmargin=5pt, innertopmargin=5pt, innerbottommargin=5pt]
    }{
        \end{mdframed}
    }
}
\makeatother

